% Rabinovich 2014, complete guarded future translation and interval complement.

\subsection{Complete-order expressiveness with a guarded past limit}
\label{sec:p09-guarded-future-completeness}

Every future monadic first-order formula has a guarded
\(U,K^-\) expression on Dedekind-complete orders. The interval
construction below supplies the negation step and the final syntactic
translation. It uses the actual past-limit operation, including its
value at an actual endpoint; it does not return to the insufficient
\(U\)-only basis. The primary is Rabinovich's Theorem~7.4 and its
supporting interval proof
\cite[secs.~5 and 7, pp.~7--15]{Rabinovich2014Kamp}.

\begin{hbdefinition}{def:p09-gfc-canonical-words}{The canonical label language and its interval words}
Let \(\mathcal A\) be the full finite \(U,K^-\) language of
\cref{def:p09-hr-guarded-past-limit}, with Boolean constants,
arbitrary Boolean operations and the specified strict until. Thus
\[
 K^-A(t)=\forall r\left(r<t\to
                 \exists u\,(r<u<t\land A(u))\right).
\]
In particular \(K^-A\) is true at a first point. A canonical
label is a supplied finite formula of \(\mathcal A\), evaluated
in the whole original model at its sole current-point argument.
Naming such a formula as a predicate is only an abbreviation;
all names are inlined by capture-free substitution in the output.

For \(n\geq1\) and a word
\(W=(A_0,B_1,A_1,\ldots,B_n,A_n)\), define
\[
\begin{split}
 [W](a,b)={}&a<b\land
 \exists z_1\cdots z_{n-1}\bigl[
      a=z_0<z_1<\cdots<z_n=b\\
 &{}\land\bigwedge_{i=0}^n A_i(z_i)
       \land\bigwedge_{i=1}^n
          \forall u\,(z_{i-1}<u<z_i\to B_i(u))\bigr].
\end{split}
\]
The displayed \(z_0,z_n\) are endpoint abbreviations, not
extra quantified variables. A one-interval word has no interior
witness. Both actual endpoints have labels. Finite word lists
denote disjunctions, with the empty list false. This is the
ordered case of the source's canonical
\((a,b)\)-decomposition grammar. Its additional cases \(a=b\)
and \(b<a\) are kept separately when complements are unguarded.
No formal infinity is ever a predicate argument.
\end{hbdefinition}

\begin{hblemma}{thm:p09-gfc-label-transfer}{proved=annex:p9.gfc.label.transfer}{Transferring the completed local construction to canonical labels}
The infimum descriptors, finite event-case compiler and bounded
positive-chain construction of
\cref{thm:p09-gpss-infimum-descriptors,thm:p09-gpss-infimum-case-compiler,thm:p09-gpss-complete-positive-chain}
remain valid when their supplied labels belong to \(\mathcal A\),
and their output is interpreted as canonical words.

The finite construction treats each distinct input label as a
temporary unary letter, applies the proved uniform construction,
then substitutes its original \(\mathcal A\) formula. Its added
Boolean and strict-future operations preserve \(\mathcal A\).
The former left label becomes \(A_0\), and its unlabelled upper
endpoint receives true. This is a substitution theorem for
these supplied constructions, not an identification of the
\(U\)-only and \(U,K^-\) languages.
\end{hblemma}

\begin{hblemma}{thm:p09-gfc-positive-assembly}{proved=annex:p9.gfc.positive.assembly}{Positive canonical assembly with all endpoint labels}
Finite canonical word lists on the same ordered endpoints
are closed under conjunction. Enumerate every order-preserving
merge of their interior chains, including equal points from
different chains. Conjoin labels at each endpoint and each
merged point, and transfer every covering gap label onto
inserted points and subgaps.

Guarded concatenation at \(a<c<b\) joins a word on \((a,c)\)
and one on \((c,b)\), conjoining their two endpoint labels
and every additional required label at \(c\). Distribute over
all finite choices. These constructions also apply when the
right operand is a future ray. An empty conjunction is the
one-interval true word, or the zero-witness true ray.
\end{hblemma}

\begin{hbdefinition}{def:p09-gfc-prefix-handler}{The supplied complement handler when the last gap is uniform}
Let \(W\) have \(n\geq2\) intervals. Define \(Q_W(a,b)\)
to say that a realization of
\[
                 (A_0,B_1,A_1,\ldots,B_{n-1},A_{n-1})
\]
starts at \(a\) and ends at a point strictly before \(b\).
Its explicit event decision is
\[
 A_0(a)\land
 V_{(A_1,\ldots,A_{n-1});(B_1,\ldots,B_{n-1})}(a,b),
\]
using \cref{def:p09-gpss-complete-positive-chain}.
Use the transferred event compiler to obtain a finite canonical
list \(\overline Q_W\) for its relative complement.

Write \(\mathsf{All}_B(a,b)\) for
\(\forall u(a<u<b\to B(u))\).
Define \(\mathsf{UC}_W(a,b)\) by the canonical list for
\[
 \begin{cases}
   \neg A_0(a)\lor\neg A_1(b),&n=1,\\
   \neg A_n(b)\lor\overline Q_W(a,b),&n\geq2.
 \end{cases}
\]
Point-only alternatives have a one-interval true gap and
true at the other endpoint.
\end{hbdefinition}

\begin{hblemma}{thm:p09-gfc-uniform-final-gap}{proved=annex:p9.gfc.uniform.final.gap}{The complete uniform-final-gap reduction}
For every canonical word \(W\), on actual \(a<b\),
\[
 \mathsf{All}_{B_n}(a,b)\ \longrightarrow\
             \bigl(\neg[W](a,b)\longleftrightarrow
                         \mathsf{UC}_W(a,b)\bigr).
\]
The handler is fully constructed from the positive-chain
event recursion and requires no interval-complement
induction hypothesis of the same length.
For \(n\geq2\) it also covers adjacent outer endpoints,
where \(Q_W\) has no possible interior witness.
\end{hblemma}

\begin{hbdefinition}{def:p09-gfc-upper-boundary}{The final-gap failure set and its actual supremum}
For \(a<b\), set
\[
                 S_{a,b}=\{u:a<u<b,\ \neg B_n(u)\}.
\]
For a proposed interior boundary \(c\), define
\[
\begin{split}
 \Sigma_W(a,c,b)={}&a<c<b\land
               \mathsf{All}_{B_n}(c,b)\\
       &{}\land\bigl(\neg B_n(c)\lor K^-(\neg B_n)(c)\bigr).
\end{split}
\]
The existential \(\exists c\,\Sigma_W(a,c,b)\) is a
canonical two-interval word: the first gap is true, the
second \(B_n\), the middle label is the displayed
disjunction, and both outer labels are true. When \(c\)
remains named, \(\Sigma_W\) is that word's matrix at
the specified middle point; the same \(c\) is retained
in every additional condition.
\end{hbdefinition}

\begin{hblemma}{thm:p09-gfc-upper-boundary}{proved=annex:p9.gfc.upper.boundary}{Every final-gap boundary case with its exact formula}
In a complete order, exactly one of the following holds:
the set \(S_{a,b}\) is empty; it is nonempty with supremum
\(b\); or it has a unique supremum \(c\) with \(a<c<b\).
The first two cases are expressed respectively by
\[
              \mathsf{All}_{B_n}(a,b),\qquad
              K^-(\neg B_n)(b).
\]
The third is expressed by \(\exists c\,\Sigma_W(a,c,b)\).
At an interior supremum the attained subcase has
\(\neg B_n(c)\); the unattained subcase has
\(B_n(c)\land K^-(\neg B_n)(c)\).
The final-gap universal in \(\Sigma_W\) is retained
in both subcases. Every past-limit label is applied
to an actual point having the earlier point \(a\),
so its witness condition is nonvacuous here.
\end{hblemma}

\begin{hbdefinition}{def:p09-gfc-word-cuts}{Every point and gap cut of a word}
For \(n\geq2\), set
\[
\begin{aligned}
 P_i&=(A_0,B_1,\ldots,B_i,A_i),&
 S_i&=(A_i,B_{i+1},\ldots,B_n,A_n)
                          &&(1\leq i<n),\\
 L_i&=(A_0,B_1,\ldots,B_{i-1},A_{i-1},B_i,B_i),&
 R_i&=(B_i,B_i,A_i,B_{i+1},\ldots,B_n,A_n)
                          &&(1\leq i\leq n).
\end{aligned}
\]
An omitted prefix or suffix has no entries; thus
\(L_1=(A_0,B_1,B_1)\) and
\(R_n=(B_n,B_n,A_n)\).
These are words, not temporal operators. Their interval
counts are \(i,n-i,i,n-i+1\), respectively.
\end{hbdefinition}

\begin{hblemma}{thm:p09-gfc-word-cuts}{proved=annex:p9.gfc.word.cuts}{The exact cut identity, including both extreme gaps}
For every actual \(a<c<b\),
\[
 [W](a,b)\longleftrightarrow
 \bigvee_{i=1}^{n-1}\bigl([P_i](a,c)\land[S_i](c,b)\bigr)
 \ \lor\
 \bigvee_{i=1}^{n}\bigl([L_i](a,c)\land[R_i](c,b)\bigr).
\]
Point cuts make \(c\) a witness and retain its point label.
Gap cuts put \(c\) inside one original open gap and therefore
require \(B_i(c)\), as well as both subgap universals.
The formula holds for each chosen interior \(c\), including
one defined by a final-gap supremum.
\end{hblemma}

\begin{hbmethod}{def:p09-gfc-complement-algorithm}{The terminating canonical complement recursion}
For a word \(W\) with \(n\geq1\) intervals, construct a list
\(\mathsf N_W(a,b)\) for its complement under \(a<b\).
At \(n=1\), return the lists for
\[
 \neg A_0(a)\ \lor\ \neg A_1(b)\
              \lor\exists u\,(a<u<b\land\neg B_1(u)).
\]
The last alternative has one interior point labelled
\(\neg B_1\), true gaps, and true endpoint labels.

For \(n\geq2\), define the three branch conditions
\[
\begin{aligned}
 E_1&=\neg A_n(b)\lor K^-(\neg B_n)(b),\\
 E_2&=A_n(b)\land\mathsf{All}_{B_n}(a,b),\\
 E_3(a,c,b)&=A_n(b)\land\neg K^-(\neg B_n)(b)
                                      \land\Sigma_W(a,c,b).
\end{aligned}
\]
Each is an explicit canonical condition. Using the words
in the cut definition, form
\[
\begin{aligned}
 D_i(a,c,b)&=\mathsf N_{P_i}(a,c)\lor\mathsf N_{S_i}(c,b)
                                      &&(1\leq i<n),\\
 H_i(a,c,b)&=\mathsf N_{L_i}(a,c)\lor\mathsf N_{R_i}(c,b)
                                      &&(2\leq i<n),\\
 H_1(a,c,b)&=\mathsf N_{L_1}(a,c)\lor\mathsf{UC}_{R_1}(c,b).
\end{aligned}
\]
Here every \(\mathsf N\) call has fewer than \(n\) intervals.
The handler \(\mathsf{UC}_{R_1}\) is the already completed
uniform-final-gap construction, not a recursive call of
length \(n\). Put
\[
 M_W(a,c,b)=
       \bigwedge_{i=1}^{n-1}D_i(a,c,b)
       \land H_1(a,c,b)
       \land\bigwedge_{i=2}^{n-1}H_i(a,c,b),
\]
where an empty conjunction is true. Return the finite
canonical expansion of
\[
 E_1\ \lor\
       (E_2\land\mathsf{UC}_W(a,b))
 \ \lor\
       \exists c\bigl(a<c<b\land E_3(a,c,b)\land M_W(a,c,b)\bigr).
\]
To expand a branch, distribute every finite list and split
any condition on \((a,b)\) at the named \(c\), using all
point and gap cuts. Merge the resulting lists separately
on \((a,c)\) and \((c,b)\), and concatenate at \(c\).
All point conditions at the three actual endpoints are
conjoined there. The proof justifies the omitted extreme
gap cut \(i=n\).
\end{hbmethod}

\begin{hbtheorem}{thm:p09-gfc-interval-complement}{proved=annex:p9.gfc.interval.complement}{Full bounded relative complement in the permitted canonical basis}
The finite algorithm above terminates and satisfies
\[
       a<b\land\neg[W](a,b)
                            \longleftrightarrow\mathsf N_W(a,b)
\]
on every Dedekind-complete model with arbitrary unary
interpretations. It retains all attained and unattained
supremum cases, witness coincidences, empty intervals,
endpoint labels and shorter-prefix conditions.

For an unguarded complement, add \(a=b\) and \(b<a\)
as separate source-order cases. Finite canonical lists
are closed under relative complement by complementing
each member and applying finite conjunction. The label
class is \(\mathcal A=TL(U,K^-)\), as required by
Rabinovich's Lemma~7.8(2), with its permitted actual
endpoint labels \cite[pp.~9--11, 14]{Rabinovich2014Kamp}.
\end{hbtheorem}

\begin{hblemma}{thm:p09-gfc-ray-translation}{proved=annex:p9.gfc.ray.translation}{The complete return from canonical future rays}
A canonical ray with actual points
\(x=z_0<z_1<\cdots<z_m\), point labels \(A_i\),
gap labels \(B_1,\ldots,B_m\), and final gap label
\(B_{m+1}\) on \((z_m,+\infty)\), has the
\(\mathcal A\) expression
\[
 T_m=A_m\land G B_{m+1},\qquad
 T_i=A_i\land(B_{i+1}\mathbin{\mathsf U_s}T_{i+1})
                              \quad(i=m-1,\ldots,0).
\]
For \(m=0\) this is \(A_0\land GB_1\).
Each permitted cofinal conjunct is translated by
\(F\Psi\land GF\Psi\). Finite ray lists translate
by disjunction; their complements are zero-witness
rays with left label equal to the negated resulting
\(\mathcal A\) formula and final gap true.
The equivalences hold on every linear order and at
last points as well.
\end{hblemma}

\begin{hblemma}{thm:p09-gfc-regional-induction}{proved=annex:p9.gfc.regional.induction}{All future regional constructors in the canonical language}
Every \(\operatorname{Bet}_G(a,b)\) component on
actual \(a<b\), and every \(\operatorname{Aft}_G(a)\)
component, has an effectively constructed finite
canonical word or ray list over complete orders.
The induction is on the exact source stage
\(\operatorname{Bet}_{G,n}\). It includes positive
and negative signed existentials, finite conjunctions
with no children, empty disjunctions, and every
permitted point literal.

Together with the complete GPSS future special normal
form, this supplies a formula \(\theta\in\mathcal A\)
for every future first-order formula. Its label names
are inlined. Only actual lower endpoints and bounded
intervals or future rays occur; no formal left ray or
unbounded-from-below sentence is assumed expressible
in this grammar.
\end{hblemma}

\begin{hbdefinition}{def:p09-gfc-initial-normalizer}{Removing unguarded past limits at the initial point}
For \(\theta\in\mathcal A\), define \(g(\theta)\)
by commuting with atoms and Boolean operations,
retaining a whole strict-until subformula unchanged,
and replacing an encountered \(K^-B\) by true.
Thus the recursion stops at each until or past-limit
node. Every surviving past-limit occurrence lies
under strict until. The true replacement is fixed
by the displayed vacuous \(K^-\) convention at a
first point.

To place an arbitrary future first-order input in
the source's strict-quantifier form, first restrict
all quantifiers to the tail \([x,+\infty)\), then
use the capture-free identities
\[
\begin{aligned}
 \exists y\geq x\,\chi&\longleftrightarrow
                \chi[x/y]\lor\exists y>x\,\chi,\\
 \forall y\geq x\,\chi&\longleftrightarrow
                \chi[x/y]\land\forall y>x\,\chi.
\end{aligned}
\]
All bound variables are first made fresh from \(x\).
This explicitly connects GPSS's inclusive future
quantifiers with the strict convention of the
positive source's Definition~7.11.
\end{hbdefinition}

\begin{hbtheorem}{thm:p09-guarded-future-completeness}{proved=annex:p9.guarded.future.completeness}{The full guarded future completeness theorem}
For every future monadic first-order formula
\(\varphi(x)\), the finite normal-form, interval,
ray and normalization constructions above produce
a guarded \(\mathcal A\) formula \(g(\theta)\)
such that
\[
     \mathcal M,x\models\varphi
                 \quad\Longleftrightarrow\quad
     \mathcal M,x\models g(\theta)
\]
for every nonempty Dedekind-complete linear-order
model and every unary interpretation.
This includes real, finite, integer and natural
orders, with no assumption of density or absence
of endpoints.

Conversely every guarded \(\mathcal A\) formula
has an explicit future first-order translation
on every linear order: expand the strict until
and past-limit clauses with fresh variables,
then restrict all quantifiers to the current
tail. Thus the guarded fragment and future
first-order logic are expressively equivalent
over complete orders.

On discrete complete orders, each outermost past-limit
subformula has only Boolean and future-until ancestors,
so it is evaluated after the current point, where an
immediate predecessor makes it false. Replacing those
entire subformulas by false, without descending into
their operands, recovers the \(U\)-only future conclusion
without an exception at an initial point.

This is the positive result with the actual
\(U,K^-\) basis. It asserts no corresponding
completeness theorem over the rationals, no
\(U\)-only completeness over the reals, and no
unrestricted almost-future completeness theorem
\cite[Theorem~7.4, pp.~13--15]{Rabinovich2014Kamp}.
\end{hbtheorem}

\begin{hbexample}{def:p09-gfc-attainment-example}{A supremum may coincide with a required witness}
On \(\mathbb R\), take \(a=0,b=2\), a two-interval
word with true endpoint labels, \(B_1=\top\),
and its sole interior point label true only at 1.
If \(B_2\) fails only at 1, the failure supremum
is attained there, yet the word is true with its
interior witness equal to that supremum: the final
gap excludes the witness.
If instead \(B_2\) fails throughout \((1/2,1)\)
and holds at 1, the supremum is unattained,
\(K^-(\neg B_2)(1)\) is true, and the same
witness still realizes the word.
If failures are cofinal below 2, no last witness
strictly below 2 can satisfy its final gap.

The guarded formula
\[
          F\bigl(P\land K^-P\bigr)
\]
expresses the existence of a future marked point
approached from below by marked points. The
past-limit occurrence is evaluated strictly after
the current point. Its first-order future form
keeps only the interval after that current point,
and conjoining \(P\) at the current point gives
the counterexample in
\cref{def:p09-hr-until-table}.
\end{hbexample}
